\documentclass{beamer}

\usepackage{graphicx}
\usepackage{amsmath}
\usepackage{amsthm}

% \theoremstyle{definition}
% \newtheorem{definition}{Definition}

% \theoremstyle{plain}
% \newtheorem{theorem}{Theorem}

\theoremstyle{remark}
\newtheorem{remark}{Remark}

\title{Week 4 CUSF State Feedback 1: Eigenvalue Assignment}
\author{Yoonjae Hwang}
\date{2 November 2025}

\begin{document}
\frame{\titlepage}

\begin{frame}
\frametitle{Motivation}
\begin{equation}
    \begin{aligned}
        \dot{x} &= Ax + Bu \\
        y &= Cx 
    \end{aligned}
\end{equation}
\begin{itemize}
    \item Previously, we only considered using the output to control our system
    \item Instability may lie in a variable that we cannot measure directly
    \item The next two weeks we will consider the controllers of the form: $u = - Kx$
\end{itemize}
\end{frame}

\begin{frame}
\frametitle{Motivation}
If $u = - Kx$, the system dynamics fundamentally change:
\begin{equation}
    \dot{x} = (A-BK)x
\end{equation}


\begin{itemize}
    \item By changing $K$, we can alter the system dynamics
    \item Handle Multiple Input Multiple Output systems (instead of cascade)
    \item Influencing the each of the variables that determine the state
    \item Question: How do we determine $K$?
\end{itemize}
\end{frame}

\begin{frame}
\frametitle{Some Linear Algebra}
If $T \in \mathcal{L}(V,W)$, the \textsl{eigenvalues} of $T$ are defined as $\lambda \in \mathbb{C}$ such that: 
\begin{equation}
    T(v) = \lambda v
\end{equation}

The corresponding \textsl{eigenvectors} of $T$ for each $\lambda$ is $v$.
If the linear map $T$ is represented by a matrix $A$, then the eigenvalues are s that satisfies:
\begin{equation}
    \det(sI - A) = 0
\end{equation}
The characteristic polynomial of $A$ is given by: 
\begin{equation}
    p(s) = det(sI - A) = s^n + c_1 s^{n-1} + \cdots + c_{n-1} s + c_n
\end{equation}
\end{frame}

\begin{frame}
\frametitle{Cayley-Hamilton Theorem}
If $p(s) = det(sI - A) = s^n + c_1 s^{n-1} + \cdots + c_{n-1} s + c_n$ is the characteristic polynomial 
of $A$, then the corresponding matrix polynomial, $p(A) = 0$ ie:
\begin{equation}
    A^n + c_1 A^{n-1} + \cdots + c_{n-1} A + c_n I = 0
\end{equation}

Consequences: we can represent $A^n$ as:
\begin{equation}
    A^n = -c_n I - c_{n-1} A - \cdots - c_1 A^{n-1}
\end{equation}
\end{frame}

\begin{frame}
\frametitle{The Matrix Exponential}
The solution to the differential equation: 
\begin{equation}
    \dot{x} = Ax
\end{equation}
is given by: 
\begin{equation}
    x(t) = e^{At}x(0)
\end{equation}
$e^{At}$ is defined by the Taylor Series of the Exponential function:
\begin{equation}
    e^{At} = I + At + \frac{A^2t^2}{2!} + \cdots
\end{equation}
\end{frame}

\begin{frame}
\frametitle{Reachability}
What set of points in the state space can be reached through the choice of a control input?
\begin{definition}[Reachability]
A linear system is reachable if for any $x_0, x_f \in \mathbb{R^n}, \exists T>0$ and $u: [0,T] \to \mathbb{R} \text{ such that } x(0) = x_0 \implies x(T) = x_f$
\end{definition}

\begin{itemize}
    \item Addresses whether it is possible to reach all points in the state space in a transiet fashion
    \item We are most concerned about the set of equilibrium points of the system
    \item $\mathcal{R}_{eq}= \{x_e: Ax_e + Bu_e = 0 \| u_e \in \mathbb{R}\}$
\end{itemize}
\end{frame}

\begin{frame}
\frametitle{Testing for Reachability}
\begin{definition}
The \textsl{reachability matrix} is given by: 
\begin{equation}
    W_r = 
    \begin{bmatrix}
        B & AB & \cdots & A^{n-1}B
    \end{bmatrix}
\end{equation}
\end{definition}
\begin{theorem}
    For single input: a linear system is reachable if and only if the reachability matrix $W_r$ is full row rank.
\end{theorem}
You can look at the proof yourself.
There are other tests for reachability one can look at as well.
\end{frame}

\begin{frame}
\frametitle{Reachable Canonical Form}
It is often convenient to change coordinates and write the dynamics of the system in transformed coordinates. 
This way we can transform unfamiliar settings into systems that we can easily analyse. 
We employ the coordinate transform $z=Tx$, to transform a system into its \textsl{reachable canonical form}. 
You can generalise this notion to multiple inputs.
\begin{equation}
\begin{aligned}
\frac{dz}{dt}
&=
\begin{pmatrix}
-a_1 & -a_2 & -a_3 & \cdots & -a_n \\
1 & 0 & 0 & \cdots & 0 \\
0 & 1 & 0 & \cdots & 0 \\
\vdots & & \ddots & \ddots & \vdots \\
0 & 0 & \cdots & 1 & 0
\end{pmatrix}
z
+
\begin{pmatrix}
1 \\ 0 \\ 0 \\ \vdots \\ 0
\end{pmatrix}
u \\
y &= 
\begin{pmatrix}
    b_1 & b_2 & \cdots & b_n
\end{pmatrix}
z + du
\end{aligned}
\end{equation}
\end{frame}

\begin{frame}
\frametitle{Reachable Canonical Form}
The characteristic polynomial for the system is given by: 
\begin{equation}
    \lambda (s) = s^n + a_1 s^{n-1} + \cdots + a_{n-1} s + a_n
\end{equation}
The reachability matrix is: 
\begin{equation}
    \tilde{W}_r = 
    \begin{pmatrix}
        \tilde{B} & \tilde{A}\tilde{B} & \cdots & \tilde{A}^{n-1} \tilde{B}
    \end{pmatrix}
    =
    \begin{pmatrix}
1 & -a_1 & a_1^2 - a_2 & \cdots & * \\
0 & 1 & -a_1 & \cdots & * \\
0 & 0 & \ddots & \ddots & \vdots \\
\vdots & & \ddots & 1 & -a_1 \\
0 & \cdots & \cdots & 0 & 1
\end{pmatrix}
\end{equation}

\begin{equation}
    \tilde{A} = TAT^{-1}, \tilde{B} = TB
\end{equation}
And the transformation is: 
\begin{equation}
    T = \tilde{W}_r W_r^{-1}
\end{equation}
\end{frame}

\begin{frame}
\frametitle{Stabilisation by State Feedback}
We consider a linear control law:
\begin{equation}
    u = -Kx + k_f r
\end{equation}
Substituting into the system: 
\begin{equation}
    \frac{dx}{dt} = (A - BK)x + B k_f r
\end{equation}
We want to determine the feedback gain K, so that the closed loop system has the characteristic polynomial: 
\begin{equation}
    p(s) = s^n + p_1 s^{n-1} + \cdots + p_{n-1} s + p_n
\end{equation}
The stability of the system is governed by the eigenvalues of $A-BK$, ie the solutions to $p(s)$.
The equilibrium point and steady-state output for the closed loop system are given by: 
\begin{equation}
\begin{aligned}
    x_e &= -(A-BK)^{-1} Bk_f r \\
    y_e &= Cx_e + Du_e
\end{aligned}
\end{equation}
\end{frame}

\begin{frame}
\frametitle{Design}
If D = 0:
\begin{equation}
    k_f = \frac{-1}{C(A-BK)^{-1}B}
\end{equation}
Using our gains K and $k_f$, we find the eigenvalues needed to design the \textsl{dynamics} of our closed loop system to satisfy our goal. 
\end{frame}

\begin{frame}
\frametitle{Eigenvalue Assignment}
Our system is: 
\begin{equation}
\begin{aligned}
    \dot{x} &= Ax + Bu, \qquad y &= Cx + Du
\end{aligned}
\end{equation}
We assume our system is reachable/controllable.
We change the coordinates by a transformation $z=Tx$, so that our system is in the reachable canonical form. 
\begin{equation}
    \dot{z} = \tilde{A} z + \tilde{B} u
\end{equation}
Our controller is given by: 
\begin{equation}
\begin{aligned}
    u &= - \tilde{K}z + k_f r = - \tilde{K} Tx + k_f r \\
      &= -\tilde{k_1} z_1 - \tilde{k_2} z_2 - \cdots - \tilde{k_n} z_n + k_f r
\end{aligned}
\end{equation}
\end{frame}

\begin{frame}
\frametitle{More}
The closed loop system then becomes: 
\begin{equation}
\begin{aligned}
\frac{dz}{dt}
&=
\begin{pmatrix}
-a_1 - \tilde{k_1} & -a_2- \tilde{k_2} & -a_3 - \tilde{k_3} & \cdots & -a_n  - k_n \\
1 & 0 & 0 & \cdots & 0 \\
0 & 1 & 0 & \cdots & 0 \\
\vdots & & \ddots & \ddots & \vdots \\
0 & 0 & \cdots & 1 & 0
\end{pmatrix}
z
+
\begin{pmatrix}
k_f \\ 0 \\ 0 \\ \vdots \\ 0
\end{pmatrix}
r \\
\end{aligned}
\end{equation}

The characteristic polynomial is then: 
\begin{equation}
    s^n + (a_1 + \tilde{k_1})s^{n-1} + \cdots + (a_{n-1} + \tilde{k}_{n-1})s + a_n + \tilde{k}_n
\end{equation}
Equating coefficients with 
the chose closed-loop polynomial we have: 

\begin{equation}
    \tilde{k}_i = p_i - a_i
\end{equation}
\end{frame}

\begin{frame}
\frametitle{The feedback gain}
In reachable canonical form:
\begin{equation}
    \tilde{K} = 
    \begin{pmatrix}
        p_1 - a_1 & p_2 - a_2 & \cdots & p_n - a_n
    \end{pmatrix}
\end{equation}
Thus overall, transforming back to its original: 
\begin{equation}
    K = \tilde{K}T = 
    \begin{pmatrix}
        p_1 - a_1 & p_2 - a_2 & \cdots & p_n - a_n
    \end{pmatrix} 
    \tilde{W}_r W_r^{-1}
\end{equation}
The feedforward gain is given by: 
\begin{equation}
    k_f = \frac{-1}{C(A-BK)^{-1}B}
\end{equation}
The method requires us to understand the dynamics of our desired system 
and to relate it to the corresponding characteristic polynomial for it. 
\end{frame}



\end{document}